\subsection*{A sublinear bound within every fixed multiplicative orbit}

The preceding bound still permits a linear-sized shared-prefix family,
which could be enough for the padding reduction. A stronger obstruction
applies within a fixed multiplicative orbit of binary subspaces.

\begin{theorem}[Transfer within a multiplicative orbit]
\label{thm:multiplicative-orbit-transfer}
Let $K=\mathbb F_{2^m}$, $N=2^m$, and let $U\subset K$ be an
$r$-dimensional binary subspace, with $1\le r<m$. For $a\in K^*$ put
\[
 H_a=\{v\in K:\operatorname{Tr}_{K/\mathbb F_2}(auv)=0
                         \text{ for all }u\in U\}.
\]
Injectively label $K\setminus\{0\}$ by elements of a field of
characteristic different from two. Let $\mathcal C$ be a collection
of distinct $H_a$ whose punctured locators share their first $s$
coefficients below the leading term, where $0\le s\le N/2^r-2$.
Define
\[
 b=\left\lceil\log_2\frac{(2^r-1)N}{s+1}\right\rceil,
 \qquad h=|\{c\in K^*:cU=U\}|,
\]
and assume $b\le m$. Then
\[
 |\mathcal C|\le\frac4h N^{\,1-2^{1-b}}.
\]
For fixed $r$ and a fixed positive prefix fraction $(s+1)/N$, this
is $o(N)$. In particular, for $r=2$, $s=N/8-1$, and $m\ge6$,
\[
 |\mathcal C|\le\frac4h N^{15/16},\qquad h\in\{1,3\}.
\]
The case $h=3$ is the family of $\mathbb F_4$-linear hyperplanes.
\end{theorem}
\begin{proof}
Let $A\subset K^*$ consist of all normals for members of
$\mathcal C$. Nondegeneracy of the trace pairing shows that
$H_a=H_{a'}$ exactly when $aU=a'U$, so $|A|=h|\mathcal C|$.
Newton identities express each power sum in the corresponding
elementary symmetric functions without division. Thus the moments
\[
 M_j(a)=\sum_{v\in H_a\setminus\{0\}}x_v^j
\]
are independent of $a\in A$ for $1\le j\le s$; the zeroth
moment is the common size $N/2^r-1$. This step does not require
the label-field characteristic to exceed $s$.

Suppose $A$ contains an affine binary $b$-flat
$a_\epsilon=a_0+\sum_{i=1}^b\epsilon_i w_i$ with independent $w_i$.
For $\beta\in\mathbb F_2^b\setminus\{0\}$ define
\[
 c_\beta(v)=2^{-b}\sum_{\epsilon\in\mathbb F_2^b}
       (-1)^{\beta\cdot\epsilon}{\bf1}_{v\in H_{a_\epsilon}}.
\]
All its moments vanish: $\sum_{v\ne0}c_\beta(v)x_v^j=0$ for
$0\le j\le s$. The character expansion of the annihilator of $U$
gives, with $\operatorname{Tr}=\operatorname{Tr}_{K/\mathbb F_2}$,
\[
 c_\beta(v)=2^{-r}
 \sum_{\substack{u\in U\setminus\{0\}\\
          \operatorname{Tr}(uw_iv)=\beta_i\ (1\le i\le b)}}
             (-1)^{\operatorname{Tr}(ua_0v)}.
\]
For fixed $u\ne0$, multiplication by $u$ preserves independence
of the $w_i$, so the displayed fiber has $N/2^b$ points. Therefore
$c_\beta$ is supported on at most $(2^r-1)N/2^b\le s+1$ points.

At least one such vector is nonzero over the label field. Otherwise
invertibility of the binary Walsh transform would make all $2^b$
incidence rows identical. A fixed $H_a$ can have at most $2^r-1$
normals: for any fixed nonzero $u_0\in U$, their distinct products
$au_0$ all belong to its same annihilator. Our choice of $b$ has
$2^b>2^r-1$, a contradiction. On the support of a nonzero $c_\beta$,
the square Vandermonde matrix of moments from zero to one less than
the support size is invertible, because the labels are distinct.
Its vanishing moments give the final contradiction. Thus $A$ has
no affine binary $b$-flat.

An elementary induction bounds an affine-$j$-flat-free set $B$ in a
binary group of size $q$ by $4q^{1-2^{1-j}}$. For $j=1$ it has at
most one point. For $j\ge2$, pair its elements along each nonzero
direction $v$. The occupied pairs form an affine-$(j-1)$-flat-free
set in the quotient of size $q/2$. With $\alpha=2^{2-j}$,
\[
 |B|(|B|-1)=\sum_{v\ne0}|B\cap(B+v)|
       \le8(q-1)(q/2)^{1-\alpha}.
\]
Consequently $|B|\le1+2\cdot2^{\alpha/2}q^{1-\alpha/2}
\le4q^{1-\alpha/2}$. Apply this with $j=b$ and $q=N$ to $A$.

Finally, $\{c\in K:cU\subseteq U\}$ is a finite subfield:
it is closed under addition and multiplication, and every nonzero
element acts bijectively on $U$. Its binary degree divides $r$.
For $r=2$ it has size two or four, so $h=1$ or $3$; the latter
case makes $U$ a scalar copy of $\mathbb F_4$.
\end{proof}

This rules out a linear-sized transfer from any bounded number of
fixed-codimension multiplicative orbits. At codimension two, a
$cN$-sized shared-prefix family would have to occupy at least
$(c/4)N^{1/16}$ orbits. Arbitrary codimension-two families can span
linearly many orbits, so the general linear-sized-transfer question
remains open.

The exact checker verifies the character support formula on all
$63$ affine binary five-flats avoiding zero at $N=64$, and $20$
such flats at $N=256$. It separately checks multiplicative orbits
with stabilizers one and three, and a codimension-three orbit with
seven-dimensional affine sections. Odd-field prefix tests, a
characteristic-two negative control, and a small complete odd-field
transfer complement these checks. They supplement the universal proof.
